\documentclass[a4paper,11pt]{article}

\usepackage[margin=0.85in]{geometry}
\usepackage[hidelinks]{hyperref}
\usepackage{xcolor}
\usepackage[T1]{fontenc}
\usepackage[utf8]{inputenc}
\usepackage{lmodern}
\usepackage{setspace}

\definecolor{accent}{HTML}{4B1FA5}
\definecolor{teal}{HTML}{157A86}
\definecolor{textgray}{HTML}{2F2F2F}

\pagestyle{empty}
\setlength{\parindent}{0pt}
\setlength{\parskip}{9pt}
\color{textgray}
\sloppy

\begin{document}

\begin{center}
  {\LARGE \textbf{Ayman Patel}}\\
  \vspace{5pt}
  {\small
  London, United Kingdom \quad | \quad
  \href{mailto:ayman.patel97@gmail.com}{ayman.patel97@gmail.com} \quad | \quad
  \href{https://linkedin.com/in/aymanapatel}{linkedin.com/in/aymanapatel} \quad | \quad
  \href{https://patelofthought.com/me}{patelofthought.com/me}
  }
\end{center}

\vspace{-3pt}
{\color{teal}\hrule height 1pt}

\vspace{12pt}
29 July 2026

\vspace{8pt}
{\color{teal}\textbf{\MakeUppercase{Senior Engineer, Servicing}}}

\vspace{5pt}
Dear Sonali,

I am applying for the Senior Engineer role in Funding Circle's Loan Servicing team. The opportunity to evolve the systems behind the full lifecycle of a loan is compelling because it combines the work I do best: building reliable backend services, making complex financial workflows easier to operate, and helping engineering teams make sound technical decisions.

During six years at Mastercard, I worked across mission-critical payment systems where accuracy, traceability, and resilience were core design requirements. I built Spring services and event-driven flows using Kafka, worked with PostgreSQL and Oracle, and introduced circuit breakers and retries across more than five microservices, improving availability by 15\%. I also used Splunk, Grafana, Dynatrace, and OpenTelemetry to reduce incident resolution time from days to hours. That experience maps closely to Funding Circle's need for servicing infrastructure that remains accurate, scalable, resilient, and easy to operate.

I would also bring the technical leadership expected of a senior engineer. I translated product requirements into technical specifications, led architecture decisions across teams in different countries, and coached more than ten engineering teams on accessible engineering patterns. My internal whitepapers became reference material shared with 200 engineers. On delivery quality, I maintained 95\% automated test coverage while helping cut an end-to-end CI suite from eight hours to three through parallel execution, selective retries, and environment caching.

My commercial backend experience is strongest in Java and Spring, while my recent Python work includes FastAPI services, data-heavy ML systems, and reliable agent workflows. At the Nebius AI Engineering Fellowship, I built Python-based systems with deterministic controls for agent outputs, distributed workloads on Kubernetes, and a shared MCP tool layer. This gives me a practical base for contributing to Funding Circle's Python services and emerging AI-powered internal tooling. I have not used Ruby professionally, and I would approach it in the same way I have successfully moved across Java, Python, Go, and infrastructure domains: by grounding learning in tests, system behaviour, and production outcomes.

Funding Circle's focus on helping UK small businesses gives this technical work a clear purpose. I would welcome the chance to bring my financial-systems background, backend engineering judgement, and team-wide approach to reliability to the Servicing team.

\vspace{6pt}
Kind regards,

\vspace{6pt}
{\color{accent}\textbf{Ayman Patel}}

\end{document}
